% =====================================================================
%  CSL（pandoc --citeproc）で組んだ書誌を LaTeX で受けるための定義。
%
%  Octavo は引用を CSL で解決するので、body.tex には biblatex の
%  \autocite ではなく「解決済みの地の文」と、末尾の CSLReferences 環境が入る。
%  その環境の定義は pandoc の standalone テンプレートが持っているものと同じ。
%  投稿先の main.tex に差し替えたら、この定義をそこに写すか、
%      octavo template copy paper/csl-preamble.tex
%  でプロジェクトの templates/ に置いて \input する。
%
%  （同梱の paper/ja/main.tex・paper/en/main.tex には最初から同じ定義が
%    書いてあるので、それを使っているならこのファイルは要らない。）
% =====================================================================
\makeatletter
\@ifundefined{CSLReferences}{%
  \newlength{\cslhangindent}
  \setlength{\cslhangindent}{1.5em}
  \newlength{\csllabelwidth}
  \setlength{\csllabelwidth}{3em}
  \newenvironment{CSLReferences}[2]% #1 ぶら下げインデント, #2 項目間の空き
   {\begin{list}{}{%
     \setlength{\itemindent}{0pt}
     \setlength{\leftmargin}{0pt}
     \setlength{\parsep}{0pt}
     \ifodd #1
      \setlength{\leftmargin}{\cslhangindent}
      \setlength{\itemindent}{-1\cslhangindent}
     \fi
     \setlength{\itemsep}{#2\baselineskip}}}
   {\end{list}}
  \newcommand{\CSLBlock}[1]{\hfill\break\parbox[t]{\linewidth}{\strut#1\strut}}
  \newcommand{\CSLLeftMargin}[1]{\parbox[t]{\csllabelwidth}{\strut#1\strut}}
  \newcommand{\CSLRightInline}[1]%
    {\parbox[t]{\dimexpr\linewidth-\csllabelwidth\relax}{\strut#1\strut}}
  \newcommand{\CSLIndent}[1]{\hspace{\cslhangindent}#1}
}{}
\makeatother
